\section{引言}

国际航运业正同时面临温室气体减排、燃料多元化和运行可靠性保障等多重压力。
国际海事组织（IMO）2023 战略及行业订单趋势表明，替代燃料推进系统已由远期概念逐步转化为工程现实
\cite{imo2023ghgstrategy,classnk2025alternativefuels}。
随着液化天然气（LNG）、甲醇和氨等燃料被用于降低碳排放强度，船用双燃料和多燃料发动机已成为大型动力系统研发中的重要方向
\cite{wingd_lng_faq_2025,wingd_ammonia_faq_2025,classnk_methanol_2024}。

然而，替代燃料的引入也给船用发动机燃烧过程带来新的挑战。
低反应性的主燃料在实际压燃条件下难以自行点火，因此发动机仍然需要高反应性且可靠的点火源。
燃料系统因替代燃料而更低碳、更灵活，但燃烧过程仍依赖一个尺度较小却十分关键的局部事件：柴油预喷
\cite{wingd_xdf_manual_2021,classnk_methanol_2024}。
对于喷油器开发而言，核心控制问题因此是预喷喷雾本身的空间发展：穿深、锥角和形态决定点火核的位置与体积，进而影响主燃料混合、火焰传播和整体燃烧效率
\cite{Karim2015DualFuel,WeiGeng2016DualFuelReview,Sahoo2009DualFuelReview}。
若预喷穿深过短，点火核可能体积不足或出现在不利区域；若穿深过长，喷雾可能偏离设计意图并增加壁面相互作用风险。
在工业研发周期中，此类预测器的价值并不是直接认证燃烧性能，而是在进入昂贵发动机试验或计算流体力学（CFD）模拟之前减少喷嘴与工况候选数量，提前标记 wall-proximity 风险，并将已有光学喷雾归档复用为低成本工程证据。

当这一控制目标未被满足时，同一次预喷还可能成为运行风险来源。
若液相包络在形成有效点火核之前到达活塞顶或缸壁，可能导致壁面润湿、润滑油稀释、混合准备恶化以及未燃碳氢排放升高
\cite{peraza2022swi,ahmad2025fueldilution,MorieiraMotaPanao2010}。
因此，实际问题在于其空间发展是否始终处在一个能够同时支持点火、混合和安全裕度的可控窗口内。

已有研究和工程实践为评估预喷发展提供了若干成熟方法：发动机试验、CFD、经验建模和光学诊断。
发动机试验能够在完整系统中考察喷雾、燃烧、壁面相互作用和排放，而高保真 CFD 则能解析更丰富的流动细节。
经验穿深关联式和现象学模型可快速描述喷雾发展，将穿深与壁距比较则构成了一阶撞壁筛查的自然几何判据
\cite{HiroyasuArai1990,NaberSiebers1996,Baumgarten2006MixFormation}。
光学喷雾实验则提供了互补路径，能够给出可重复的液相包络观测；其中，高速 Mie 成像尤其适合观察液相边界和宏观穿深演化
\cite{Linne2013SprayImaging}。

然而，上述方法都只能覆盖快速、条件参数化、实验依据明确的预喷穿深预测需求中的一部分。
发动机试验成本高、周期长；CFD 需要复杂输入、模型调参与显著计算资源。
经验模型结构紧凑且便于解释，但短预喷会使其使用更加困难，因为可观测窗口可能主要由喷射起始与喷射终止瞬态占据
\cite{ZhouLiYi2021,ZhouLiLaiWei2019}。
光学实验能够提供图像证据，但图像后处理结果仍是离散观测表，而不是可直接用于快速设计迭代的预测模型。

因此，本文要解决的是位于高通量光学喷雾归档与更高成本船用发动机验证之间的一个具体中间层缺口。
这一缺口由三方面约束共同定义：多孔喷嘴会产生名义上相似但物理上不同的喷油束；有限视场会使部分晚期穿深历史成为右删失观测；原始 Mie 强度又会受到照明、眩光和多重散射影响。
如果把每条观测轨迹的可见末端直接视为真实终点，模型就可能把“离开图像”误学为“物理饱和”。
因此，所缺的不是新的喷雾物理理论，也不是额外的单点诊断测量，而是一个位于光学实验归档与计算或实验成本高昂的验证之间、快速、可重复、具备不确定度意识且范围明确的前置筛查层。

数据驱动代理模型为这一目标提供了可行方向。
近期研究表明，神经网络可以根据工况参数预测穿深和锥角等喷雾宏观特性
\cite{VazKarathanasis2026,liu2020sprayreview}。
在本文的设置下，同一光学配置下获得的多批次 Mie 记录可以转化为逐喷油束时序观测量，并用于训练从工况到穿深轨迹的映射。
因此，本文旨在开发并评估一条可审计的工作流，把大规模 Mie 散射喷雾视频归档转化为位于光学实验与高成本验证之间的早期工程预测层。
该工作流被有意定位为 design triage 工具：它根据 Mie 定义穿深和几何 wall-proximity 风险对候选预喷喷雾几何进行排序，而点火质量、燃烧效率、排放和有害壁面润湿仍留给后续发动机代表性验证。
论文的目标包括：
\begin{enumerate}
  \item 从高速非反应 Mie 成像视频中提取逐喷油束几何观测量；
  \item 基于有限视场观测构造穿深轨迹目标；
  \item 训练能够同时预测穿深均值和条件不确定性的代理模型；
  \item 将预测分布连接到简化代表性发动机几何，形成概率性撞壁筛查代理指标。
\end{enumerate}

本文的主要贡献包括一套 CUDA 加速的视频到观测量工作流、一套用于重跑大型光学 campaign 的可复现归档与 checkpoint 约定、一套面向右删失穿深目标的降阶拟合策略、一套多阶段不确定性感知学习框架，以及一种简化的概率性 wall-proximity 筛查方法。
本文范围有意限定于受控加压舱条件下非反应、非蒸发的柴油预喷喷雾。
该模型应被理解为其数据支持域内的工程筛查工具，而非任意未知工况下均可靠的通用预测器、直接发动机燃烧预测器，也不是已经完成验证的真实液态撞壁真值模型。

第 2 章从光学、计算机视觉、并行编程、优化和机器学习等多个知识领域出发，建立文献依据并提供本文所需的理论基础；
第 3 章描述实验数据与图像处理工作流；
第 4 章描述轨迹目标、代理模型建模和概率性筛查；
第 5 章报告实验结果；
第 6 章讨论哪些内容已经得到验证以及当前适用边界；
第 7 章区分已验证结果、原型边界与下一步工作。
